% ===================================================================
%  TAVOLA N. 15 — Il mercato. --- I generi alimentari.
%  Traduction 2026 d'apres texto/io, controlee sur texto/fr et
%  texto/en. Consigne : voir texto/it/10-tavola-01.tex.
%
%  DEUX MOTS DE CETTE PAGE PARTAGENT L'EUROPE EN DEUX, ET LA LIGNE
%  PASSE PAR LA POLOGNE.
%
%  LA TOMATE DE L'ALINEA 3 (14). Le fruit arrive du Mexique au XVIe
%  siecle avec son nom nahuatl, « tomatl », et l'Europe entiere le
%  prend : tomate, tomato, Tomate, tomat, domates. L'italien ne l'a
%  pas pris. Il a regarde le fruit et l'a appele pomodoro,
%  LA POMME D'OR — parce que les premieres varietes qu'on a vues
%  etaient jaunes.
%
%  Or ce nom-la est reparti vers l'est, et il y est reste :
%     polonais    « pomidor »       ukrainien  « помідор »
%     russe       « помидор »       lituanien  « pomidoras »
%     bielorusse  « памідор »       bulgare    « домат » (turc)
%  UNE MOITIE DE L'EUROPE APPELLE LA TOMATE PAR SON NOM AZTEQUE ET
%  L'AUTRE PAR SA COULEUR, et la frontiere entre les deux passe a peu
%  pres par la Vistule. Aucune des deux moities ne sait qu'elle a
%  choisi.
%
%  LA POMME DE TERRE DE L'ALINEA 10 (61), ET C'EST UN ECHANGE
%  D'ERREURS ENTRE VOISINS. Le tableau 15 allemand avait montre que
%  « Kartoffel » est l'italien « tartufolo », LA TRUFFE — l'Allemagne
%  a pris a l'Italie un mot qui designe un champignon. Et l'italien,
%  lui, dit patata, qui est l'espagnol « patata », lequel est
%  le taino « batata », LA PATATE DOUCE, croise avec le quechua
%  « papa ». L'Italie a donc pris a l'Espagne un mot qui designe une
%  autre plante.
%
%     CHACUN DES DEUX VOISINS A PRIS A L'AUTRE OU A SON VOISIN UN NOM
%     QUI NE CONVENAIT PAS, ET LES DEUX ERREURS TIENNENT ENCORE.
%
%  Le tableau 15 estonien avait la meme plante et la meme hesitation,
%  « kartul » contre « maaõun » ; le tableau 15 allemand avait
%  Kartoffel, Erdapfel et Grundbirne. Quatre langues, cinq noms, et
%  pas un qui decrive la chose.
%
%  ENFIN, CE MARCHE-CI EST LE PLUS DIVISE DES SEIZE, plus encore que
%  l'allemand. L'allemand avait douze couples nord-sud sur cette page ;
%  l'italien en aurait trente si l'on descendait au-dessous de la
%  langue commune — anguria et cocomero, melone et popone, cetriolo et
%  citrolo, prezzemolo et petrosino, pane et micca. On ecrit la langue
%  commune, celle des dictionnaires et de la television ; les vingt
%  autres noms de la pasteque ne sont pas des fautes, ce sont d'autres
%  langues, et elles n'ont pas de colonne ici.
% ===================================================================

%%K t15-apar-1 apar
\VUsaut{4.0mm}
\VUtitre{15.0pt}{60}{TAVOLA N. 15}
\VUsaut{3.0mm}

%%K t15-tit-1 sub
\VUpk{11.4pt}{40}{Il mercato. --- I generi alimentari.}
\VUsaut{3.0mm}

%%K t15-01-1 p
1. --- La maggior parte dei commercianti che ci forniscono i generi alimentari
necessari è riunita sotto la \VUgras{tettoia} \textsuperscript{(1)} del
\VUgras{mercato coperto}, o nelle vie vicine.

%%K t15-02-1 p
2. --- Il \VUgras{salumiere} \textsuperscript{(2)}, che vende \VUgras{prosciutto}
\textsuperscript{(3)}, \VUgras{sanguinaccio} \textsuperscript{(4)},
\VUgras{salsiccia} \textsuperscript{(5)}, \VUgras{salsiccia di trippa}
\textsuperscript{(6)} e \VUgras{lardo} \textsuperscript{(7)}, sta accanto alla
\VUgras{venditrice di burro e formaggio} \textsuperscript{(8)}, il cui banco è
apparecchiato con \VUgras{formaggi olandesi} \textsuperscript{(9)} dalla crosta rossa,
\VUgras{camembert} \textsuperscript{(10)}, grandi \VUgras{forme di groviera}
\textsuperscript{(11)} e \VUgras{panetti di burro} \textsuperscript{(12)} fresco.

%%K t15-03-1 p
3. --- Davanti a lei una \VUgras{verduraia} \textsuperscript{(13)} offre alle sue
clienti bei \VUgras{pomodori} \textsuperscript{(14)}, \VUgras{cavolfiori}
\textsuperscript{(15)}, \VUgras{lattuga} \textsuperscript{(16)}, \VUgras{cavoli}
\textsuperscript{(17)}, \VUgras{zucche} \textsuperscript{(19)}, \VUgras{meloni}
\textsuperscript{(18)} e perfino \VUgras{funghi} \textsuperscript{(20)}. Ma un
\VUgras{garzone del birraio}, che ha fermato il suo \VUgras{carro}
\textsuperscript{(21)} carico di \VUgras{botti di birra} \textsuperscript{(22)} proprio
davanti al suo banco, impedisce alle clienti di avvicinarsi. Un ortolano le porta un
cesto di \VUgras{carciofi} \textsuperscript{(23)}.

%%K t15-tit-2 sub
\VUsaut{4.0mm}
\VUpk{10.2pt}{40}{Vendere e comprare.}
\VUsaut{3.0mm}

%%K t15-04-1 p
4. --- Al mercato c'è un gran movimento, perché è l'ora dell'\VUgras{asta}
\textsuperscript{(24)}. Le \VUgras{massaie} \textsuperscript{(26)} che vengono qui a
fare la spesa si fermano lungo la strada davanti al banco della
\VUgras{trippaia} \textsuperscript{(25)} per comprare la \VUgras{trippa} o una
\VUgras{testina di vitello}.

%%K t15-05-1 p
5. --- Una \VUgras{cuoca} \textsuperscript{(27)} robusta contratta un bellissimo
\VUgras{tonno} con la \VUgras{pescivendola} \textsuperscript{(28)}, che alza i prezzi
come le pare. Le è arrivata una grossa partita di \VUgras{anguille},
\VUgras{sogliole}, \VUgras{trote} e \VUgras{carpe}. Il suo \VUgras{merluzzo} e le sue
\VUgras{cozze} sono freschissimi.

%%K t15-tit-3 sub
\VUsaut{4.0mm}
\VUpk{10.2pt}{40}{Roba a buon mercato e roba cara.}
\VUsaut{3.0mm}

%%K t15-06-1 p
6. --- La \VUgras{pollivendola} \textsuperscript{(29)}, che sta accanto a lei, è molto
onesta. Quando vende un \VUgras{tacchino}, un'\VUgras{oca}, un'\VUgras{anatra} o un
semplice \VUgras{pollo}, non chiede più del giusto prezzo.

%%K t15-07-1 p
7. --- All'angolo della piazza, al primo piano, si è messa da poco una
\VUgras{negoziante di biancheria} \textsuperscript{(30)}, dove i compratori trovano
sempre un'ottima scelta di \VUgras{tovaglie}, \VUgras{tovaglioli},
\VUgras{grembiuli} e \VUgras{strofinacci}. Sullo stesso piano un
\VUgras{coltellinaio} \textsuperscript{(31)} ha messo la sua bottega. Affila i
\VUgras{coltelli} delle donne del mercato e fa agli ortolani \VUgras{roncole}
del \VUgras{acciaio} più duro.

%%K t15-tit-4 sub
\VUsaut{4.0mm}
\VUpk{10.2pt}{40}{La drogheria.}
\VUsaut{3.0mm}

%%K t15-08-1 p
8. --- Sotto la sua bottega ha aperto da poco un grande negozio di alimentari. Non è
più la semplice, modesta \VUgras{drogheria} \textsuperscript{(32)} di una volta, che
vendeva solo generi di uso quotidiano --- \VUgras{tè} \textsuperscript{(33)},
\VUgras{cioccolato} \textsuperscript{(34)}, \VUgras{zucchero in pani}
\textsuperscript{(37)} e \VUgras{sapone} \textsuperscript{(38)}. Qui si vende di tutto:
\VUgras{biscotti secchi}, \VUgras{conserve} \textsuperscript{(35)},
\VUgras{liquori} \textsuperscript{(36)}, \VUgras{rum}, \VUgras{acquavite},
\VUgras{kirsch}, \VUgras{anisetta} e \VUgras{curaçao}. La selvaggina si trova qui con
la stessa facilità --- \VUgras{lepre} \textsuperscript{(39)}, \VUgras{tordi}
\textsuperscript{(40)}, \VUgras{pernici} \textsuperscript{(41)}, \VUgras{quaglie}
\textsuperscript{(42)}, \VUgras{anatre selvatiche} \textsuperscript{(43)} o
\VUgras{fagiano} \textsuperscript{(44)} --- come la frutta tropicale, per esempio le
\VUgras{banane} \textsuperscript{(46)} e gli \VUgras{ananas}.

%%K t15-09-1 p
9. --- Un \VUgras{commesso} \textsuperscript{(45)} molto gentile del negozio è pronto a
dare tutto quello che gli si chiede: \VUgras{marmellata} \textsuperscript{(47)}, di
ribes o di cotogne, \VUgras{strutto} \textsuperscript{(48)},
\VUgras{cetriolini} \textsuperscript{(49)} o \VUgras{sardine} \textsuperscript{(50)}
salate.

%%K t15-09-2 p
Questo è molto comodo per i \VUgras{clienti} \textsuperscript{(51)}, che accanto alla
merce più comune trovano anche le primizie più rare e la frutta più buona: succose
\VUgras{pere} \textsuperscript{(52)}, \VUgras{prugne} \textsuperscript{(53)},
\VUgras{uva} \textsuperscript{(54)}, \VUgras{pesche} \textsuperscript{(55)},
\VUgras{mandorle} \textsuperscript{(56)} e \VUgras{noci} \textsuperscript{(57)}.

%%K t15-tit-5 sub
\VUsaut{4.0mm}
\VUpk{10.2pt}{40}{I clienti.}
\VUsaut{3.0mm}

%%K t15-10-1 p
10. --- Il \VUgras{riso} \textsuperscript{(58)} e le \VUgras{lenticchie}
\textsuperscript{(59)} stanno accanto alla cassa delle \VUgras{aringhe}
\textsuperscript{(60)} affumicate; la \VUgras{patata} \textsuperscript{(61)} e il
\VUgras{fagiolo} \textsuperscript{(63)} stanno vicino alle \VUgras{mele}
\textsuperscript{(62)}, ai \VUgras{fichi} \textsuperscript{(64)}, alle
\VUgras{prugne secche} \textsuperscript{(65)} e alle \VUgras{nocciole}
\textsuperscript{(66)}. Questo negozio vende anche \VUgras{uova}
\textsuperscript{(67)} fresche e \VUgras{olive} \textsuperscript{(68)}, e così i
\VUgras{cesti} \textsuperscript{(70)} e le \VUgras{sporte a rete}
\textsuperscript{(73)} si riempiono in fretta.

%%K t15-11-1 p
11. --- Il \VUgras{caffè} \textsuperscript{(72)} di questa ottima ditta si è fatto un
nome meritato fra le \VUgras{domestiche} \textsuperscript{(69)} e le cuoche, perché il
vecchio \VUgras{droghiere} \textsuperscript{(71)} lo tosta lui stesso con la massima
cura.

%%K t15-tit-6 sub
\VUsaut{4.0mm}
\VUpk{10.2pt}{40}{Cose da vedere.}
\VUsaut{3.0mm}

%%K t15-12-1 p
12. --- Non tutti vengono al mercato per comprare da mangiare. Nel pittoresco andirivieni
della viuzza che porta alla tettoia si vedono le persone più diverse. Accanto a un
bonario \VUgras{macellatore} \textsuperscript{(75)} che spinge davanti a sé una
\VUgras{carriola} \textsuperscript{(74)}, ecco due soldati in licenza, molto fieri di
mostrare le loro divise sgargianti: l'\VUgras{ussaro} \textsuperscript{(76)} con il suo
\VUgras{dolman} azzurro e il \VUgras{colbacco piumato}, e il
\VUgras{caporale degli zuavi} con i calzoni larghi e il \VUgras{fez} in testa.

%%K t15-13-1 p
13. --- Il \VUgras{fornaio} \textsuperscript{(80)}, in piedi sulla soglia del
\VUgras{forno} \textsuperscript{(78)}, ha appena impastato il suo \VUgras{pane}
\textsuperscript{(79)} e tolto dal forno i \VUgras{cornetti} \textsuperscript{(81)}, i
\VUgras{panini} \textsuperscript{(82)} e le \VUgras{pagnotte} \textsuperscript{(83)}
che sono in vetrina.

%%K t15-14-1 p
14. --- Sopra il forno abita una brava \VUgras{sarta} \textsuperscript{(84)} che dà
lavoro a diverse \VUgras{ricamatrici} \textsuperscript{(85)}. Accanto a lei abita una
\VUgras{lavandaia e stiratrice} \textsuperscript{(86)}, che inamida la
\VUgras{biancheria} \textsuperscript{(88)} dopo averla lavata e asciugata, e la stira
con un \VUgras{ferro} \textsuperscript{(87)}.

%%K t15-tit-7 sub
\VUsaut{4.0mm}
\VUpk{10.2pt}{40}{Carne, pesce e verdura.}
\VUsaut{3.0mm}

%%K t15-15-1 p
15. --- Nella \VUgras{macelleria} \textsuperscript{(89)} al pianterreno si vende carne
di ogni genere --- \VUgras{montone} \textsuperscript{(90)}, \VUgras{manzo}
\textsuperscript{(94)} o \VUgras{vitello} \textsuperscript{(92)} --- come
\VUgras{cosciotti}, \VUgras{bistecche}, \VUgras{arrosti} e \VUgras{costolette}. Il
\VUgras{garzone del macellaio} \textsuperscript{(93)} taglia tutta questa carne e mette
da parte gli \VUgras{ossibuchi} \textsuperscript{(96)}. Sulla porta della macelleria sta
una \VUgras{venditrice di ostriche} \textsuperscript{(97)} che vende \VUgras{ostriche}
\textsuperscript{(98)}. Se uno vuole, gliele apre.

%%K t15-16-1 p
16. --- Tutti questi commercianti pagano una licenza o un posteggio; per questo la
\VUgras{guardia} \textsuperscript{(100)} non ha nessuna pietà per le povere
\VUgras{verduraie ambulanti} \textsuperscript{(99)}, che fanno loro concorrenza girando
con i loro carretti carichi di \VUgras{carote}, \VUgras{ravanelli},
\VUgras{rape}, \VUgras{porri} o mazzi di \VUgras{asparagi}, \VUgras{piselli
freschi}, \VUgras{spinaci}, \VUgras{acetosa}, \VUgras{crescione} e
\VUgras{prezzemolo}.

%%K t15-tit-8 sub
\VUsaut{4.0mm}
\VUpk{10.2pt}{40}{Attenti alla guardia!}
\VUsaut{3.0mm}

%%K t15-17-1 p
17. --- Il ragazzo che vende \VUgras{aglio} \textsuperscript{(101)} e \VUgras{cipolle}
sa benissimo che neanche lui può vendere la sua roba vicino al mercato. Scappa via, con
il rischio di rovesciare il cesto di \VUgras{granchi}, \VUgras{gamberi di fiume} e
\VUgras{gamberetti} della \VUgras{venditrice} \textsuperscript{(102)} di
\VUgras{aragoste} e di \VUgras{astici}.

%%K t15-18-1 p
18. --- Sono tutti in movimento e fanno un gran baccano; ma si sente lo stesso
chiaramente il richiamo del \VUgras{vetraio} \textsuperscript{(103)} ambulante e il
grido del \VUgras{carbonaio} \textsuperscript{(104)}, che ha appena lasciato per un
momento il suo carro per consegnare un \VUgras{sacco di carbone}
\textsuperscript{(105)}.
\VUsaut{6.0mm}
\VUfilet{22mm}
